\subsection{International Governance and Diverging National Rules}
\label{sec:8.7.8}
Almost none of the international governance of AI is binding. What has emerged instead is mostly a set of voluntary, overlapping, specialized bodies, closer to how international environmental and trade norms developed piecemeal than to a single founding charter — with one recent exception.

The Council of Europe's Framework Convention on Artificial Intelligence and Human Rights, Democracy and the Rule of Law is the first legally binding international treaty specifically on AI. Negotiated by the Council of Europe's member states alongside outside participants that included the European Union and the United States, it requires five ratifications, at least three of them Council of Europe members, before it enters into force \autocite{coe2024framework}. What it binds signatories to is written at the level of a human-rights treaty and not a technical code: consistency with human rights, democracy, and the rule of law across an AI system's lifecycle, with the operational detail of how a signatory meets that obligation left to domestic law. It is a foundation binding AI law can be built on, and not a substitute for the domestic statutes discussed below.

Everything else in this space is voluntary, and the roll is longer than it is deep. The OECD's Recommendation on Artificial Intelligence was the first intergovernmental AI standard, and its principles became the basis for the G20's own \autocite{oecd2019recommendation}. The European Commission's High-Level Expert Group on AI produced “Ethics Guidelines for Trustworthy AI,” whose seven requirements later shaped the structure of the EU AI Act \autocite{aihleg2019ethics}. UNESCO's Recommendation on the Ethics of Artificial Intelligence, adopted by all 193 member states, is the broadest formal endorsement any AI ethics instrument has received \autocite{unesco2021recommendation}. The Global Partnership on AI is a standing venue for research and policy coordination between governments \autocite{gpai2020joint}, and the Partnership on AI does the same work across companies and researchers \autocite{deepmind2016partnership}. The G7's Hiroshima AI Process pairs guiding principles for every actor in the AI lifecycle with a code of conduct for organizations building the most advanced systems, and carries legal force in none of the participating countries \autocite{g72023hiroshima}. A narrower, more technical layer sits alongside all of it: the Bletchley Declaration committed 29 signatories to cooperate specifically on frontier-model risk, and the national AI Safety Institutes it produced later joined an International Network whose members agreed to share testing methodology and run joint evaluations of the same frontier models instead of each working from a separate standard \autocite{aisafetysummit2023bletchley} \autocite{nist2024network}.

The summit series that produced the Bletchley Declaration continued past it and, in doing so, showed how quickly this kind of voluntary consensus can fray. A follow-up summit in Seoul in 2024 kept the same safety framing; a third, in Paris in February 2025, rebranded itself the AI Action Summit and shifted its emphasis from risk toward adoption and opportunity, and its closing statement was not signed by the United States or the United Kingdom, two of the series' own founding participants \autocite{aiactionsummit2025paris}. A coordination mechanism that depends on every major AI developer's home government agreeing to the same framing is only as durable as that agreement, and within two years of Bletchley, the agreement had already come apart.

None of these bodies can compel a state or a company that ignores it. Each is a standard, a set of principles, or a forum, useful for building the shared vocabulary that eventually makes binding law possible, and that gap between voluntary standard and binding law is the default state of international coordination on anything individual states have reasons to handle differently. The EU AI Act, discussed below, is the first instance of this kind of groundwork becoming enforceable law, and it took most of a decade of standard-setting exactly like this to get there. The proliferation carries its own cost in the meantime: a company or a smaller government trying to comply with “international AI governance” faces a shifting handful of overlapping standards, each with different membership and scope and no formal mechanism for reconciling a conflict between them, which is why “the international community agrees” is, at the level of any specific rule, usually false. What states have done once left to write their own rules is the subject of what follows.

\runin{Democracy, Authoritarianism, and Diverging National Rules}
Section~\ref{sec:6.4.1} names six documented forms of authoritarian AI misuse, among them social scoring, mass biometric surveillance, and automated censorship, and section~\ref{sec:6.4.3} lays out the policy levers available against them: export controls on surveillance-capable systems, laws that name specific prohibited designs rather than regulating AI in the abstract, and developer-facing standards that move faster than legislation does. What belongs here is not a second pass over that ground but what individual states have done with the room those levers leave open, since geopolitics is the study of how different governments use the same available tools differently.

The European Union chose comprehensive statute. The AI Act sorts systems by risk tier, a small set of practices banned outright under Article 5, a larger tier of “high-risk” applications subject to conformity assessment and ongoing monitoring, and most other uses left largely unregulated, rather than writing one rule meant to cover every case \autocite{eu2024aiact}. Section~\ref{sec:6.4.3} already covers Article 5's specific prohibitions; the structural point worth adding here is that the EU decided the risk a system poses, not the technology it uses, should determine how much scrutiny it gets. Enforcement runs through a dedicated AI Office inside the European Commission, which can demand technical documentation from a general-purpose model's developer, run its own evaluations, and fine a noncompliant provider a percentage of global turnover — reaching any company that offers a model inside the EU regardless of where the company is based, the same extraterritorial mechanism that let the GDPR set a de facto floor for privacy law well beyond Europe's borders \autocite{eu2024aiactenforcement}.

The United States made the opposite structural bet: no comprehensive federal AI statute, and enforcement distributed instead across existing sectoral regulators, the Federal Trade Commission, the Equal Employment Opportunity Commission, and the Food and Drug Administration among them, applying laws written before AI existed to harms AI now causes. The National Institute of Standards and Technology's voluntary AI Risk Management Framework serves as a common technical reference point across those agencies rather than a binding rule of its own \autocite{nist2023airmf}. The wager is that existing consumer-protection, employment, and safety law already reaches most AI harms, and that naming AI specifically risks regulating the label rather than the conduct.

The United Kingdom, after leaving the EU's regulatory orbit, made a third choice: rather than create a new AI regulator or a new statute, its 2023 white paper directed existing sector regulators to apply five shared principles, safety, transparency, fairness, accountability, and contestability, within their existing remits \autocite{dsit2023proinnovation}. The wager is closer to the American one than the European: that domain expertise inside an existing regulator beats a new institution built from scratch, even at the cost of less uniform treatment across sectors.

China took a fourth path, one that does not fit neatly next to the other three because it starts from a different premise about what regulation is for. Its Cyberspace Administration, working with several other agencies, issued rules governing generative AI services in 2023 that require providers to file their algorithms with the state before public release, disclose the source and scale of their training data on request, and ensure that generated content reflects the government's designated social values \autocite{cac2023generativeai}. The EU, US, and UK approaches all regulate AI as a private product that might harm a consumer, a worker, or a subject of an automated decision, and argue with each other only about how centralized that regulation should be. China's rules regulate AI as a communications channel the state has a standing interest in controlling content on, which is a different kind of law aimed at a different kind of risk, and comparing it to the other three on the same axis, more regulation or less, misses what actually distinguishes it.

Among the EU, US, and UK, there is an argument about which structural bet is safer, and it runs in both directions. A comprehensive statute can freeze a category of harm into law before anyone has seen what a fast-moving technology actually does with it, and lock in a definition of “high-risk” that the next generation of systems slides around rather than falls into. A sectoral approach leaves gaps precisely where no existing regulator has jurisdiction, since a new capability does not always map onto an old agency's mandate. Neither failure mode is hypothetical, and none of the three statutes has existed long enough to show which one it falls into more often.

The EU, US, and UK differ on how centralized AI regulation should be, but not on what it is for: constraining what a private company can build and sell. None of the three, including the EU's binding statute, does anything to stop a government from building the authoritarian applications section~\ref{sec:6.4.1} documents within its own borders, because a law that governs private companies is a different instrument from one that constrains what a government does with AI systems it controls directly. China's rules show what a law aimed at the second target actually looks like — and that it can coexist with, rather than substitute for, exactly the kind of state control section~\ref{sec:6.4.1} catalogs. Conflating the two kinds of law, treating “AI is regulated here” as an answer to “is AI used against this country's own population here,” is a recurring error worth naming rather than repeating.
